\section{Сепарабельные замыкания}

Пусть $\Om$~--- поле, содержащее $F$, и $\mathcal E$~--- непустое множество промежуточных полей $F\subset E\subset\Om$ со
свойством
\[
(*)\qquad\text{для любых }E_1,E_2\in\mathcal E\text{ найдётся }E\in\mathcal E\text{ с }E_1,E_2\subset E.
\]
Тогда $E(\mathcal E)=\bigcup_{E\in\mathcal E}E$~--- подполе $\Om$ (\emph{направленное объединение}): по $(*)$ любое
конечное множество элементов $E(\mathcal E)$ лежит в одном $E\in\mathcal E$, а потому там же лежат их суммы,
произведения и обратные.

\begin{lemma}[6.8b: транзитивность сепарабельности]
Пусть $F\subset E\subset K$~--- алгебраические расширения, $K/E$ и $E/F$ сепарабельны. Тогда $K/F$ сепарабельно.
\end{lemma}

\begin{proof}
Пусть $\alpha\in K$ и $g=X^n+e_{n-1}X^{n-1}+\dots+e_0$~--- его минимальный многочлен над $E$; он сепарабелен. Поле
$E_0=F(e_0,\dots,e_{n-1})$ конечно над $F$ и порождено сепарабельными элементами, поэтому сепарабельно над $F$ (\rf{3.15}).
Многочлен $g$ неприводим над $E\supset E_0$, тем более над $E_0$, так что $\alpha$ сепарабелен над $E_0$. В
характеристике $0$ доказывать нечего; пусть $\operatorname{char}F=p>0$. Сначала отметим:

\emph{если $\beta$ сепарабелен над полем $L$ и $\beta^p\in L$, то $\beta\in L$:} минимальный многочлен $\beta$ над $L$
делит $X^p-\beta^p=(X-\beta)^p$, т.\,е.\ равен $(X-\beta)^r$, а сепарабельность даёт $r=1$.

Элемент $\alpha$ сепарабелен над $E_0(\alpha^p)\supset E_0$, поэтому $\alpha\in E_0(\alpha^p)$. Поле
$E_0(\alpha^p)$ порождено над $F(\alpha^p)$ элементами $E_0$, которые сепарабельны над $F$, а значит, и над
$F(\alpha^p)$; по \rf{3.15} $E_0(\alpha^p)/F(\alpha^p)$ сепарабельно. Итак, $\alpha$ сепарабелен над $F(\alpha^p)$ и
$\alpha^p\in F(\alpha^p)$, откуда $\alpha\in F(\alpha^p)$, т.\,е.\ $F(\alpha)=F(\alpha^p)$.

Пусть теперь минимальный многочлен $\mu$ элемента $\alpha$ над $F$ несепарабелен. Тогда $\mu=\mu_1(X^p)$ (\rf{2.20}),
и $\mu_1$~--- минимальный многочлен $\alpha^p$ (он неприводим, иначе разлагался бы и $\mu$). Получаем
$[F(\alpha):F]=\deg\mu=p\deg\mu_1=p\,[F(\alpha^p):F]$, что противоречит равенству $F(\alpha)=F(\alpha^p)$. Значит,
$\alpha$ сепарабелен над $F$.
\end{proof}

Применим направленные объединения к множеству подполей $E\subset\Om$, конечных и сепарабельных над $F$. Композит двух
таких подполей снова конечен и сепарабелен над $F$ (он порождён сепарабельными элементами, \rf{3.15}), так что их
объединение $L$~--- подполе $\Om$. Поле $L$ сепарабельно над $F$, и каждый элемент $\Om$, сепарабельный над $F$, лежит в
$L$ (он порождает конечное сепарабельное подрасширение); иначе говоря, $L$~--- сепарабельное замыкание $F$ в $\Om$
(\rf{3.15}). Если $\Om$ алгебраично над $F$, то оно \emph{чисто несепарабельно} над $L$: если $\alpha\in\Om$ сепарабелен над
$L$, то $L(\alpha)/L$ и $L/F$ сепарабельны, по \rf{6.8b} $\alpha$ сепарабелен над $F$, т.\,е.\ $\alpha\in L$.

\begin{definition}[6.9]
\begin{enumerate}[label=(\alph*)]
\item Поле $\Om$ \emph{сепарабельно замкнуто}, если каждый непостоянный сепарабельный многочлен из $\Om[X]$
раскладывается в $\Om$ на линейные множители.
\item Поле $\Om$~--- \emph{сепарабельное замыкание} своего подполя $F$, если $\Om$ сепарабельно и алгебраично над $F$ и
сепарабельно замкнуто.
\end{enumerate}
\end{definition}

\begin{theorem}[6.10]
\begin{enumerate}[label=(\alph*)]
\item У каждого поля есть сепарабельное замыкание.
\item Пусть $E$~--- сепарабельное алгебраическое расширение $F$, а $\Om$~--- сепарабельное замыкание $F$. Тогда
существует $F$-гомоморфизм $E\to\Om$; если $E$ тоже сепарабельное замыкание $F$, то каждый такой гомоморфизм~---
изоморфизм.
\end{enumerate}
\end{theorem}

\begin{proof}
Можно повторить доказательства для алгебраических замыканий, заменив слово <<многочлен>> словами <<сепарабельный
многочлен>>. Проще вывести теорему из уже доказанного.

(a) Пусть $\bar F$~--- алгебраическое замыкание $F$ и $\Fsep$~--- сепарабельное замыкание $F$ в $\bar F$ (см.~выше).
Оно сепарабельно и алгебраично над $F$. Пусть $g\in\Fsep[X]$~--- непостоянный сепарабельный многочлен. Каждый его
корень $\beta\in\bar F$ сепарабелен над $\Fsep$ (минимальный многочлен $\beta$ делит $g$), а значит, по \rf{6.8b} и над $F$,
т.\,е.\ $\beta\in\Fsep$. Так как $g$ раскладывается в $\bar F$, он раскладывается и в $\Fsep$.

(b) Пусть $\bar\Om$~--- алгебраическое замыкание $\Om$; оно алгебраично над $F$ (\rf{1.31}) и потому является
алгебраическим замыканием $F$. По \rf{6.6} есть $F$-гомоморфизм $\varphi\colon E\to\bar\Om$. Для $x\in E$ элемент
$\varphi(x)$~--- корень минимального многочлена $x$ над $F$, который сепарабелен; поэтому $\varphi(x)$ сепарабелен над
$F$, а значит, и над $\Om$. Минимальный многочлен $\varphi(x)$ над $\Om$ сепарабелен и неприводим, а $\Om$
сепарабельно замкнуто, так что его степень равна $1$, т.\,е.\ $\varphi(x)\in\Om$. Итак, $\varphi\colon E\to\Om$.

Если $E$ тоже сепарабельное замыкание, то поле $\varphi(E)\cong E$ сепарабельно замкнуто, а $\Om$ сепарабельно над
$F$ и тем более над $\varphi(E)\supset F$. Для $y\in\Om$ минимальный многочлен $y$ над $\varphi(E)$ сепарабелен и
неприводим, поэтому имеет степень $1$, т.\,е.\ $y\in\varphi(E)$. Значит, $\varphi$ сюръективен, а инъективен он как
гомоморфизм полей.
\end{proof}

\begin{center}
\begin{tikzpicture}
\node (F) at (0,0) {$F$};
\node at (0.85,0) {$\subset$};
\node (S) at (1.85,0) {$\Fsep$};
\node at (2.85,0) {$\subset$};
\node (A) at (3.75,0) {$\bar F$};
\draw[decorate,decoration={brace,mirror,amplitude=5pt},cold] (-0.25,-0.3) -- node[below=6pt,font=\footnotesize]{сепарабельно} (2.3,-0.3);
\draw[decorate,decoration={brace,mirror,amplitude=5pt},cnew] (2.4,-0.3) -- node[below=6pt,font=\footnotesize,align=center]{чисто\\ несепарабельно} (4.05,-0.3);
\end{tikzpicture}
\end{center}
\figcap{Сепарабельное замыкание внутри алгебраического. Если $F$ совершенно, то $\Fsep=\bar F$ (ср.~\rf{7.6a}); группа
$\Gal(\Fsep/F)$~--- абсолютная группа Галуа поля $F$ (гл.~7).}

\begin{aside}[6.10a: аксиома выбора в этой главе]
Сведём вместе, что где использовалось.
\begin{center}
\begin{tabular}{p{7.2cm}p{6.6cm}}
\toprule
Утверждение & Что использует \\
\midrule
Замыкание конечного или счётного поля (введение) & ничего: выборы делаются по нумерации \\
\rf{6.3} (максимальные идеалы) & лемма Цорна \\
\rf{6.4e}--\rf{6.4h} (мощности) & аксиома выбора (\rf{6.4f} ей равносильно) \\
Первое и третье доказательства существования & простой идеал~--- принцип ультрафильтра \\
Второе доказательство существования & лемма Цорна и мощности \\
\rf{6.6} (вложение в замыкание) & лемма Цорна; в счётном случае~--- выбор корней из конечных множеств (без выбора недоказуемо уже для $F=\Q$, иначе любое алгебраическое замыкание $\Q$ было бы изоморфно $\Qal$ вопреки \rf{6.4a}) \\
\rf{6.6} (<<каждый гомоморфизм~--- изоморфизм>>) & ничего \\
\rf{6.7} (изоморфизм замыканий) & принцип ультрафильтра (теорема Тихонова для конечных дискретных пространств) \\
Компактность групп Галуа (\rf{7.8}) & то же, что в \rf{6.7} \\
\bottomrule
\end{tabular}
\end{center}
Без какой-либо формы выбора ни существование, ни единственность алгебраического замыкания в общем случае не доказать
(\rf{6.4a}). Для конечных и счётных полей и подполей $\CC$ выбор не нужен для \emph{существования}: замыкание строится
явно; единственность же без выбора недоказуема уже для $\Q$ (\rf{6.4a}).
\end{aside}
